
\subsubsection{4.1 Research Questions}

We design experiments to answer the following questions:

\textbf{RQ1 (Existence):} Does SSL-SGD training create severe worst-group accuracy degradation on spurious correlation benchmarks, relative to supervised training? \textit{(H-E1 gate --- confirmed)}

\textbf{RQ2 (Geometric Diagnosis):} Does the sharpness anisotropy ratio of converged SSL models exceed 1.2, and does it correlate negatively with worst-group accuracy across training checkpoints? \textit{(H-E1 anisotropy test --- pending)}

\textbf{RQ3 (Geometric Intervention):} Does replacing SGD with SAM during SSL pretraining reduce the anisotropy ratio and improve worst-group accuracy, without group annotations? \textit{(H-M2 intervention test --- planned)}

Each RQ corresponds to a mechanism step in the causal chain: SSL-SGD creates anisotropy (RQ1/RQ2), SAM reduces it (RQ3). RQ1 is confirmed by current results. RQ2 and RQ3 are complete experimental protocols awaiting full validation from 200-epoch runs.

\subsubsection{4.2 Datasets}

\textbf{Waterbirds} \citep{Sagawa2019} is our primary benchmark. It superimposes CUB-200-2011 bird images onto Places365 backgrounds with 95\% spurious correlation: 95\% of waterbirds appear on water backgrounds and 95\% of landbirds appear on land backgrounds at training time. The four demographic groups are (waterbird, water), (waterbird, land), (landbird, water), (landbird, land), with training sizes 3498, 184, 56, 1057 respectively. Validation and test sets are balanced across groups. We use the standard splits from kohpangwei/group\_DRO.

\textbf{CelebA} \citep{liu2015celeba} (planned) correlates hair color (blonde/non-blonde) with gender. We use WILDS \citep{koh2021wilds} for downloading. CelebA was excluded from the initial experiment due to a WILDS server HTTP 500 error; results are deferred to an extended evaluation.

\textbf{CMNIST} \citep{arjovsky2019invariant} (planned) correlates digit color with binary label at 99\% strength. Excluded from the initial fast run due to speed constraints; included in the full 9-combination evaluation protocol.

\subsubsection{4.3 Models and Baselines}

\textbf{SSL Methods:}
\begin{itemize}
\item \textbf{DINO} \citep{caron2021dino}: Self-distillation SSL using ViT-S/8 backbone (384-dim features). Primary model in H-E1 validation. DINO uses publicly released ImageNet-pretrained ViT-S/8 weights rather than training from scratch on Waterbirds; see Limitation 0 in Section 6.2 for implications.
\item \textbf{SimCLR} \citep{chen2020simclr}: NT-Xent contrastive loss on ResNet-50 (2048-dim). Trained from scratch on Waterbirds. Linear probe evaluated at 10 epochs (fast run); 200-epoch full run ongoing.
\item \textbf{MoCo-v2} \citep{chen2020simclr}: InfoNCE with momentum encoder. Included in full 9-combination protocol.
\end{itemize}

\textbf{Baselines:}
\begin{itemize}
\item \textbf{SGD-trained SSL} (primary): Standard SSL pretraining with SGD optimizer, followed by linear probe evaluation.
\item \textbf{Supervised} (oracle upper bound): Supervised cross-entropy training (ViT-S/16) with identical evaluation protocol. WGA 81.93\% on Waterbirds.
\item \textbf{LFR} \citep{ghaznavi2023lfr}: Annotation-free post-hoc resampling applied to SSL features.
\item \textbf{EVaLS} \citep{ghaznavi2024evals}: Extended LFR with environment inference.
\item \textbf{Group DRO} \citep{Sagawa2019}: Oracle baseline requiring group labels (84.6\% WGA on Waterbirds).
\end{itemize}

\subsubsection{4.4 Implementation Details}

\textbf{SSL pretraining:} ResNet-50 backbone with identity final layer (2048-dim features). DINO uses ViT-S/8 with pre-trained weights from the official DINO repository (ImageNet-pretrained; see Limitation 0). Standard augmentation: RandomResizedCrop(224), RandomHorizontalFlip, ColorJitter(0.4, 0.4, 0.4, 0.1), RandomGrayscale(p=0.2), GaussianBlur. Batch size 256, cosine learning rate schedule, 200 pretraining epochs.

\textbf{SAM configuration:} rho=0.05 using davda54/sam wrapper [davda54, 2021]; ASAM variant (rho=0.5, adaptive) included as secondary variant.

\textbf{Linear probe evaluation:} SGD optimizer, lr=0.01, 100 epochs, no group labels used. Worst-group accuracy computed over 4 groups using kohpangwei/group\_DRO protocol.

\textbf{Anisotropy measurement:} N=100 random direction samples, probe epochs=100, SAM rho=0.05, checkpoint epochs {50, 100, 150, 200}.

\textbf{Hardware:} Experiments ran on a GPU cluster (CUDA 12.9, PyTorch). cuDNN was disabled due to a torch+cu124 / CUDA 12.9 driver mismatch; this slows training but does not affect numerical results.

\textbf{Reproducibility:} Seed 1 for initial PoC run. Code built on izmailovpavel/spurious\_feature\_learning and kohpangwei/group\_DRO frameworks.

\subsubsection{4.5 Evaluation Metrics}

\textbf{Primary:}
\begin{itemize}
\item \textbf{Worst-group accuracy (WGA):} Minimum accuracy across 4 demographic groups on the test set. Evaluated using kohpangwei/group\_DRO protocol. Reported as percentage.
\item \textbf{Sharpness anisotropy ratio (AR):} $\text{AR}(\theta) = \Delta\mathcal{L}_{\text{spur}} / \Delta\mathcal{L}_{\text{rand}}$. Gate threshold: AR > 1.2.
\item \textbf{Pearson correlation (r):} Correlation between AR and WGA across 4 checkpoints (epochs 50/100/150/200). As noted in Section 3.3, $n=4$ requires $|r|>0.95$ for $p<0.05$; results will be reported with this constraint acknowledged.
\end{itemize}

\textbf{Secondary:}
\begin{itemize}
\item \textbf{Average accuracy:} Mean accuracy across all groups. Monitors core feature preservation under SAM.
\item \textbf{Linear probe proxy precision/recall:} Precision and recall of top-25\% high-loss samples against held-out group labels. Required validation of the annotation-free proxy in SSL representations (see Section 3.3 Assumption and Limitation 3).
\end{itemize}

